\label{chap:conclusion}

BioSeqDB demonstrates that sequence search can be represented as a PostgreSQL
execution operation rather than as a scalar call to an external mapper. The
system stores packed biological sequences, constructs persistent minimizer
indexes, and runs a minimap2-derived pipeline for seeding, chaining, filtering,
and final extension. Planner support and custom scans preserve the batch
structure required by SIMD and GPU implementations while presenting typed SQL
results that can participate in filters, projections, and joins.

The first research question concerned a shared representation for the
algorithmic stages. Alphabet-generic compact ranges cover resident values,
while lazy sequence views represent fragmented, sliced, and composed sources.
Algorithms read the required region without first converting the complete
value to a byte array. The second question concerned heterogeneous execution. Minimizer
extraction and chaining provide regular work for the implemented SIMD and GPU
routes. Dispatch follows explicit algorithm settings, input compatibility,
index form, and resource availability. Final alignment uses banded or unbanded
KSW2, adaptive Block Aligner, and GPU routes that may return specialized scripts
before device dynamic programming. Optional prefilters may reject unsuitable
candidates but do not replace alignment.

The third question concerned database execution. A custom scan owns the search
expression, batches input rows, and substitutes match fields into planned
expressions. This design also permits multiple alignments for one input row.
The fourth question concerned the surrounding data system. Physical sequence
storage, persistent indexes, bounded ingestion, external-provider relation definitions,
cache coordination, and source-driven SQL generation give the kernels explicit
storage and lifecycle contracts.

The remaining limits are operational and algorithmic. Some remote retrieval
paths still require record-sized or archive-sized memory. Provider quotas are
not coordinated across all PostgreSQL backends. Hardware-specific routes have
batch-size thresholds and require broader evaluation on additional devices.
GPU chaining may examine a bounded predecessor set, and seed and alignment
backends retain route-specific contracts. The present mapper is centered on
linear references rather than pangenome graphs. A relational database can
nevertheless plan and execute batched
sequence analysis while retaining typed data, bounded resources, persistent
derived state, and explicit algorithm selection.
